\section{Related work}
\label{sec:related}

\subsection{Bottom-up grouping and its coupling to the detector}
\label{sec:related_grouping}

% --- from chapter2_grouping.tex 2.2.1-2.2.4 ---
Most of the bottom-up grouping literature falls into four lineages.
In the first, OpenPose~\citep{cao2017openpose} predicts a part-affinity field for every limb alongside the keypoint
heatmaps, scores candidate joint pairs by the line integral of the field,
and assembles the skeletons one limb at a time by greedy bipartite matching;
PifPaf~\citep{kreiss2019pifpaf} refines the fields but keeps the greedy assembly.
In the second, associative embedding~\citep{newell2017associative} attaches a head that predicts a tag per keypoint,
each tag being pulled towards its person's mean and pushed away from other people's means,
and grouping is done by clustering the tags;
HigherHRNet~\citep{cheng2020higherhrnet} applies the mechanism to the high-resolution HRNet backbone~\citep{sun2019hrnet}.
In the third, PersonLab~\citep{papandreou2018personlab} regresses offsets between keypoints of the same person,
DEKR~\citep{geng2021dekr} regresses offsets from a person centre to its keypoints,
and CenterGroup~\citep{braso2021centergroup} associates keypoints with centres through transformer attention.
The fourth is set prediction: PETR~\citep{shi2022petr}, ED-Pose~\citep{yang2023edpose},
GroupPose~\citep{liu2023grouppose} and QueryPose~\citep{xiao2022querypose} follow DETR~\citep{carion2020detr}
and let each pose query predict a complete skeleton, with no separate grouping stage.
In all four lineages the grouping representation is learned jointly with the detector's features.
This joint training lets the grouping adapt to the detector's errors
and explains much of the accuracy these methods reach.
It also means that the grouping cannot be moved to another detector:
the fields, tags and offsets are tied to the backbone's resolution, stride and features,
and set prediction exposes no per-keypoint representation at all.

\subsection{Graph-based grouping and graph networks on pose}
\label{sec:related_graph}

% --- from chapter2_grouping.tex 2.2.5 and chapter2_gnn_pose.tex 2.3.1-2.3.4 ---
A smaller body of work approaches grouping as a graph problem.
Hierarchical graph grouping (HGG)~\citep{jin2020hgg} merges sub-graphs of detected keypoints with a learned discriminator,
but the graph and the merge predictor are trained jointly with the detector.
Yang et al.~\citep{yang2023instanceaware} supervise the self-attention of a vision transformer to encode person membership,
which places the grouping inside the detector backbone.
Graph edge classification (GEC)~\citep{maher2025gec} is the closest precedent to this work:
it operates on the output of an existing detector, predicts a same-person label for pairs of keypoints
and recovers people as connected components.
It is detector-independent, but it produces a binary decision per edge rather than a representation per keypoint.

Elsewhere in pose research, graph neural networks are used widely.
Skeleton-based action recognition builds spatial-temporal graphs over the joints of a single person~\citep{yan2018stgcn,shi2019twostreamadaptive},
and 3D lifting applies graph convolutions to one person's 2D skeleton~\citep{zhao2019semanticgcn,ci2019optimizing,liu2020gastnet}.
In both cases the keypoints are already grouped,
and the output is a label or coordinates, not an identity signal per keypoint.
The graph and clustering literature offers learned alternatives to a plain read-out:
deep embedded clustering~\citep{xie2016dec}, slot attention~\citep{locatello2020slotattention},
supervised graph partitioning from face clustering~\citep{yang2019face,yang2020gcnve},
and modularity maximisation~\citep{newman2006modularity} with its differentiable relaxation DMoN~\citep{tsitsulin2023dmon}.
This work evaluates these alternatives to the $k$-means read-out for keypoint grouping.

\subsection{Position of this work}
\label{sec:related_gap}

% --- from chapter2_gap.tex 2.6.1 ---
Detectors are now strong enough that grouping is the main remaining source of error in bottom-up systems,
yet every established grouping mechanism is tied to its detector.
Modular grouping on detached keypoints has a single detector-independent precedent,
and it classifies edges.
Graph networks on pose, by contrast, have been used almost entirely on keypoints that are already grouped.
No prior work in this modular setting trains a detector-independent graph network
that produces a per-keypoint embedding for multi-person grouping,
or estimates the number of people within the same pipeline.